\documentclass[11pt]{article}
\usepackage[a4paper,margin=0.85in]{geometry}
\usepackage{amsmath,amssymb,mathtools}
\usepackage{booktabs}
\usepackage{enumitem}
\usepackage{xcolor}
\usepackage{hyperref}
\hypersetup{hidelinks}

\newcommand{\wt}{\widetilde}
\newcommand{\Fcal}{\mathcal{F}}
\newcommand{\dd}{\mathrm{d}}

\title{Direct-Route Check of the Updated Shortcut Calculation}
\author{}
\date{7 June 2026}

\begin{document}
\maketitle

\section*{Purpose and conclusion}

This note follows the original time-convolution route in the updated
manuscript and performs the calculation explicitly.  The result is a corrected
direct-route version of the long Eq. (41) expansion, together with a check showing
that the same expression recompresses to the closed form
\begin{equation}
  \boxed{
  \wt F_\rho(z)=
  \frac{aT_{L-\rho}(y)+U_{\rho-1}(y)+U_{L-\rho-1}(y)}
  {aT_L(y)+U_{L-1}(y)}}
  \label{eq:main-closed}
\end{equation}
where
\[
  N=2L,\qquad
  a=\frac{q}{\beta(1-q)},\qquad
  y=1+\frac{1/z-1}{q},\qquad U_{-1}=0.
\]
Thus the closed form is not a different model or a shortcut around the
calculation.  It is the simplified form of the original convolution route after
all intermediate poles have been summed.

\section{Corrected starting point}

Let
\[
  \lambda=\beta(1-q).
\]
When the shortcut is \(u\to v\) and \(v\) is the absorbing target, the shortcut
is killing in the transient state space.  The corrected first-passage
generating function is
\begin{equation}
  \boxed{
  \wt F_{n_0}(v,z)=
  \wt{\Fcal}_{n_0}(v,z)
  +qz\,\wt W_{n_0}(u,z)
  \left[1-\wt{\Fcal}_{u}(v,z)\right]\wt H(z)}
  \label{eq:F-correct-H}
\end{equation}
with
\begin{equation}
  \wt H(z)=
  \frac{\beta(1-q)}{q}\,
  \frac{1}{1+z\beta(1-q)\wt W_u(u,z)}.
  \label{eq:H-def}
\end{equation}
For the antipodal case \(|u-v|=N/2\), this is
\begin{equation}
  \wt H(z)=\frac{1}{a+\phi(z)}.
  \label{eq:H-a-phi}
\end{equation}
The sign in \eqref{eq:F-correct-H} is positive.  Therefore the corrected
time-domain convolution is
\begin{equation}
  \boxed{
  F_{n_0}(v,t)=\Fcal_{n_0}(v,t)
  +q\sum_{t'=0}^{t-1}W_{n_0}(u,t-1-t')
  \sum_{t''=0}^{t'}
  \left[\delta_{t',t''}-\Fcal_u(v,t'-t'')\right]H(t'')}
  \label{eq:time-convolution}
\end{equation}
for \(t\ge1\).  This is the corrected version of the updated manuscript's
Eq. (40).

\section{Mode expansions used in the direct route}

For \(N=2L\), the no-shortcut first-passage probability is
\begin{equation}
  \Fcal_{n_0}(v,t)=
  \sum_{\ell=1}^{L}f_\ell(n_0,v)\alpha_\ell^{t-1},
  \qquad t\ge1.
  \label{eq:F0-time}
\end{equation}
The killed propagator is
\begin{equation}
  W_{n_0}(u,t)
  =
  \sum_{r=1}^{N-1}G^{(1)}_r(n_0,u,v)\gamma_r^t
  +
  \sum_{r=1}^{L}g^{(2)}_r(n_0,u,v)\alpha_r^t.
  \label{eq:W-time}
\end{equation}
Here
\[
  \alpha_r=1-q+q\cos\frac{(2r-1)\pi}{N},
  \qquad
  \gamma_r=1-q+q\cos\frac{2\pi r}{N}.
\]

The corrected \(H\)-poles are the roots of the corrected denominator.  Write
\[
  s=1/z,\qquad y_s=1+\frac{s-1}{q}.
\]
For the antipodal case
\begin{equation}
  \wt H(z)=
  \frac{T_L(y)}{aT_L(y)+U_{L-1}(y)}.
  \label{eq:H-T-over-D}
\end{equation}
Define
\[
  D(y)=aT_L(y)+U_{L-1}(y).
\]
Let
\[
  h_0=H(0)=\frac{\beta(1-q)}{q}=\frac1a.
\]
Since \(\wt H(z)\) has numerator and denominator of the same degree, its
partial-fraction expansion has a constant term.  Equivalently, \(H(0)\) must be
kept separate.  If \(D(y_k)=0\) and \(s_k=1-q+qy_k\), then for \(t\ge1\)
\begin{equation}
  H(t)=\sum_k c_k s_k^{t-1},
  \qquad
  \boxed{c_k=\frac{qT_L(y_k)}{D'(y_k)}}.
  \label{eq:H-residue}
\end{equation}
This residue formula is the safest way to write the corrected version of
Eq. (39): it is tied directly to the corrected denominator \(D(y)\).  In
practice one should solve \(D(y)=0\) itself, not only the trigonometric
parametrization \(y=\cos\theta\).  Some roots can have \(y<-1\) while still
giving an admissible \(|s_k|<1\) pole after \(s_k=1-q+qy_k\).

\section{The convolution kernels}

Equation \eqref{eq:time-convolution} contains two kinds of sums, and each must
be split into the \(H(0)=h_0\) contribution and the positive-time \(H(c)\),
\(c\ge1\), pole contribution.

\subsection*{Two-factor convolution}

The \(H(0)\) part gives simply
\begin{equation}
  A^{(0)}_t(x)=x^{t-1}.
  \label{eq:A0-kernel}
\end{equation}
For a positive-time pole mode \(s^{c-1}\), with \(a+c=t-1\) and \(c\ge1\),
\begin{align}
  \sum_{\substack{a+c=t-1\\c\ge1}}x^a s^{c-1}
  &=\sum_{c=1}^{t-1}x^{t-1-c}s^{c-1}\notag\\
  &=\boxed{\frac{s^{t-1}-x^{t-1}}{s-x}}.
  \label{eq:A-kernel-derivation}
\end{align}
Define
\begin{equation}
  \boxed{A^{(+)}_t(x;s)=\frac{s^{t-1}-x^{t-1}}{s-x}}.
  \label{eq:A-kernel}
\end{equation}
If \(x=s\), the expression is understood as the limit
\[
  A^{(+)}_t(s;s)=(t-1)s^{t-2}.
\]
For \(t=1\), \(A^{(+)}_1=0\), as there is no room for a positive-time
\(H\)-factor.

\subsection*{Three-factor convolution}

The negative part of \(1-\Fcal_u(v,\cdot)\) contributes a first-passage mode
\(y^{b-1}\) with \(b\ge1\).  The \(H(0)\) contribution is
\begin{equation}
  \boxed{
  B^{(0)}_t(x,y)=
  \sum_{\substack{a+b=t-1\\b\ge1}}x^a y^{b-1}
  =
  \frac{x^{t-1}-y^{t-1}}{x-y}.}
  \label{eq:B0-kernel}
\end{equation}
For the positive-time \(H\)-pole contribution, \(a+b+c=t-1\), \(b\ge1\),
and \(c\ge1\):
\begin{equation}
  \sum_{\substack{a+b+c=t-1\\b\ge1,\,c\ge1}}x^a y^{b-1}s^{c-1}
  =
  [z^{t-3}]
  \frac{1}{(1-zx)(1-zy)(1-zs)}.
\end{equation}
Define
\begin{equation}
  \boxed{
  B^{(+)}_t(x,y;s)=
  [z^{t-3}]
  \frac{1}{(1-zx)(1-zy)(1-zs)}.}
  \label{eq:B-kernel}
\end{equation}
When \(x,y,s\) are distinct,
\begin{equation}
  \boxed{
  B^{(+)}_t(x,y;s)=
  \frac{x^{t-1}}{(x-y)(x-s)}
  +\frac{y^{t-1}}{(y-x)(y-s)}
  +\frac{s^{t-1}}{(s-x)(s-y)}.}
  \label{eq:B-kernel-distinct}
\end{equation}
Coincident cases are obtained by taking limits.  This is where terms such as
\(t\alpha^t\) can appear in an uncompressed expansion.  Such terms are not
automatically true tail terms; they may cancel after the complete expression is
recombined.

\section{Corrected direct-route Eq. (41)}

Substitute \eqref{eq:F0-time}, \eqref{eq:W-time}, and \eqref{eq:H-residue} into
\eqref{eq:time-convolution}, keeping the \(H(0)\) part separate.  The corrected
long expression is
\begin{equation}
\boxed{\begin{aligned}
F_{n_0}(v,t)
&=
\sum_{\ell=1}^{L}f_\ell(n_0,v)\alpha_\ell^{t-1}
\\
&\quad
+q h_0
\left[
\sum_{r=1}^{N-1}G^{(1)}_r(n_0,u,v)\gamma_r^{t-1}
+
\sum_{r=1}^{L}g^{(2)}_r(n_0,u,v)\alpha_r^{t-1}
\right]
\\
&\quad
+q\sum_k c_k
\left[
\sum_{r=1}^{N-1}G^{(1)}_r(n_0,u,v)A^{(+)}_t(\gamma_r;s_k)
+
\sum_{r=1}^{L}g^{(2)}_r(n_0,u,v)A^{(+)}_t(\alpha_r;s_k)
\right]
\\
&\quad
-q h_0\sum_{\ell=1}^{L}f_\ell(u,v)
\left[
\sum_{r=1}^{N-1}G^{(1)}_r(n_0,u,v)B^{(0)}_t(\gamma_r,\alpha_\ell)
\right.
\\
&\hspace{38mm}\left.
+
\sum_{r=1}^{L}g^{(2)}_r(n_0,u,v)B^{(0)}_t(\alpha_r,\alpha_\ell)
\right]
\\
&\quad
-q\sum_k c_k\sum_{\ell=1}^{L}f_\ell(u,v)
\left[
\sum_{r=1}^{N-1}G^{(1)}_r(n_0,u,v)B^{(+)}_t(\gamma_r,\alpha_\ell;s_k)
\right.
\\
&\hspace{38mm}\left.
+
\sum_{r=1}^{L}g^{(2)}_r(n_0,u,v)B^{(+)}_t(\alpha_r,\alpha_\ell;s_k)
\right].
\end{aligned}}
\label{eq:hard-route-result}
\end{equation}
This is the corrected direct-route replacement for the updated manuscript's
Eq. (41).  Compared with the printed expression, the baseline term must target
\(v\), the global sign inherits the positive sign in Eq. (40), and the
\(g^{(2)}\) modes are paired with \(\alpha_r\), not with \(\gamma_r\).

\section{Why the direct route recompresses to the closed form}

Now specialize to \(N=2L\), \(v=0\), \(u=L\), and let \(\rho\) be the folded
distance from the starting site to the absorbing target.  The ingredients in
\eqref{eq:F-correct-H} have compact Chebyshev forms:
\begin{align}
  \wt{\Fcal}_\rho(0,z)&=\frac{T_{L-\rho}(y)}{T_L(y)},
  \label{eq:F0-rho-compact}\\
  \wt W_\rho(L,z)&=\frac{U_{\rho-1}(y)}{zq\,T_L(y)},
  \label{eq:WrhoL-compact}\\
  1-\wt{\Fcal}_{L}(0,z)&=1-\frac1{T_L(y)}
  =\frac{T_L(y)-1}{T_L(y)},
  \label{eq:one-minus-FL}\\
  \wt H(z)&=\frac{T_L(y)}{aT_L(y)+U_{L-1}(y)}.
  \label{eq:H-compact-again}
\end{align}
Substitution into \eqref{eq:F-correct-H} gives
\begin{align}
  \wt F_\rho(z)
  &=
  \frac{T_{L-\rho}}{T_L}
  +
  qz\frac{U_{\rho-1}}{zqT_L}
  \frac{T_L-1}{T_L}
  \frac{T_L}{aT_L+U_{L-1}}\notag\\
  &=
  \frac{T_{L-\rho}(aT_L+U_{L-1})+U_{\rho-1}(T_L-1)}
  {T_L(aT_L+U_{L-1})}.
  \label{eq:pre-closed}
\end{align}
The only remaining algebra is a Chebyshev identity:
\begin{equation}
  T_L(y)U_{L-\rho-1}(y)
  =
  T_{L-\rho}(y)U_{L-1}(y)-U_{\rho-1}(y).
  \label{eq:closing-identity}
\end{equation}
Using \eqref{eq:closing-identity} in \eqref{eq:pre-closed} gives exactly
\eqref{eq:main-closed}:
\[
  \wt F_\rho(z)=
  \frac{aT_{L-\rho}(y)+U_{\rho-1}(y)+U_{L-\rho-1}(y)}
  {aT_L(y)+U_{L-1}(y)}.
\]

This proves that the closed form is the direct-route expression after the
intermediate \(\alpha,\gamma,s\) terms have been recombined.

\section{Tail decay from the recompressed expression}

Let
\[
  N_\rho(y)=aT_{L-\rho}(y)+U_{\rho-1}(y)+U_{L-\rho-1}(y),
  \qquad
  D(y)=aT_L(y)+U_{L-1}(y).
\]
The poles of \(\wt F_\rho(z)\) are the roots \(D(y_k)=0\).  With
\(s_k=1-q+qy_k\), simple roots give
\begin{equation}
  \boxed{
  F_\rho(t)=\sum_k B_{\rho k}s_k^{t-1},
  \qquad
  B_{\rho k}=\frac{qN_\rho(y_k)}{D'(y_k)}.}
  \label{eq:closed-time}
\end{equation}
The long-time tail is therefore
\begin{equation}
  \boxed{
  F_\rho(t)\sim B_{\rho1}s_1^{t-1},
  \qquad
  \frac{F_\rho(t+1)}{F_\rho(t)}\to s_1.}
  \label{eq:tail}
\end{equation}
It is not generically \(t\alpha_1^t\).

\section{Check of the figure in the updated PDF}

The figure in the updated PDF is not completely useless.  Interpreted narrowly,
it is a graphical representation of the poles of \(H\): the intersections of
\[
  \phi(1/s)
\]
with the corrected horizontal line
\[
  -a=-\frac{q}{\beta(1-q)}.
\]
That part is consistent with the corrected denominator \(a+\phi(1/s)=0\).

However, the figure is better treated as a pole-location visualization rather
than as a standalone way to infer a finite transition at
\[
  s_1=\gamma_1=1-q+q\cos\frac{2\pi}{N}.
\]
For the corrected equation,
\[
  D(\cos\theta)=a\cos(L\theta)+\frac{\sin(L\theta)}{\sin\theta}=0,
  \qquad
  s_1=1-q+q\cos\theta_1,
\]
and the largest root satisfies
\begin{equation}
  \boxed{\gamma_1<s_1<\alpha_1}
  \qquad (0<\beta\le1).
  \label{eq:s1-between}
\end{equation}
At \(s=\gamma_1\), equivalently \(\theta=\pi/L\), the \(\phi\)-term is zero, so
the corrected denominator is \(a>0\), not zero.  Thus \(s_1=\gamma_1\) is only
approached in the infinite-shortcut-strength limit \(a\to0\), not reached for
finite positive admissible \(\beta\).

There is also a caption typo in the updated PDF: the sentence saying that
\(q>1\) is required for a nonzero shortcut should be replaced.  The shortcut
mass is \(\beta(1-q)\), so one needs \(0<q<1\) and \(\beta>0\), not \(q>1\).

\subsection*{Suggested replacement caption sentence}

\begin{quote}
The intersections with the horizontal line
\(-q/[\beta(1-q)]\) give the poles \(s_k\) of the corrected auxiliary
function \(\wt H_{N,N/2}(z)=1/(q/[\beta(1-q)]+\phi(z))\).  This figure is a
visualization of the \(H\)-poles only; it does not by itself determine the
finite-time one-peak/two-peak transition.  For finite \(0<\beta\le1\), the
dominant corrected pole satisfies \(\gamma_1<s_1<\alpha_1\), so the proposed
finite crossing \(s_1=\gamma_1\) does not occur.
\end{quote}

\section{Validation summary}

\begin{center}
\small
\begin{tabular}{@{}p{0.66\linewidth}p{0.26\linewidth}@{}}
\toprule
Check & Result \\
\midrule
Corrected plus/plus first-passage formula vs finite stochastic matrix
& \(2.08\times10^{-17}\) max difference \\
Corrected direct-route Eq. \eqref{eq:hard-route-result} vs finite stochastic matrix
& \(1.32\times10^{-16}\) max difference \\
Updated Eq. (32) denominator sign vs finite stochastic matrix
& \(3.29\times10^{-3}\) max difference \\
Updated Eq. (33) minus prefactor vs finite stochastic matrix
& \(6.34\times10^{-2}\) max difference \\
Closed form \eqref{eq:main-closed} vs finite stochastic matrix
& \(7.63\times10^{-17}\) max difference \\
\bottomrule
\end{tabular}
\normalsize
\end{center}

For \(N=100\), \(q=2/3\), and the fixed shortcut \(u=6\to v=56\), the clear
visual double-peak scan finds clear double peaks up to about \(\beta=0.03\) for
\(n_0=1,2,3\), and no clear double peak at
\(\beta_c=2q/[(1-q)N]=0.04\).  This supports treating bimodality as a
finite-time shape question, not as a direct consequence of the asymptotic tail
crossing.

\section*{Deliverable recommendation}

One useful package to send back is this direct-route note together with the
short benchmark note.  The direct-route note shows how the corrected long
convolution is obtained and why it recompresses to the closed form; the
benchmark note gives the numerical evidence for updating the remaining signs in
the displayed first-passage formulae.

\end{document}
